% Source-derived chapter 03: Bialynicki-Birula theory on the moduli space of Higgs bundles
% Original source: very-stable-higgs-bundles-equivariant-multiplicity-mirror-symmetry
\chapter{Bialynicki-Birula theory on the moduli space of Higgs bundles}
\label{very-stable-higgs-bundles-equivariant-multiplicity-mirror-symmetry-chapter-03}
\source{very-stable-higgs-bundles-equivariant-multiplicity-mirror-symmetry}

\begin{remark}[Source section]
\label{very-stable-higgs-bundles-equivariant-multiplicity-mirror-symmetry-chapter-03-source}
\source{very-stable-higgs-bundles-equivariant-multiplicity-mirror-symmetry}
\source{very-stable-higgs-bundles-equivariant-multiplicity-mirror-symmetry}
This chapter reproduces the corresponding section of the retrieved
LaTeX source, including its definitions, statements, examples, and
proofs. It has not yet been translated into Lean.
\notready
\end{remark}

% The source section title was: Bialynicki-Birula theory on the moduli space of Higgs bundles
\label{bbhiggs}
\source{very-stable-higgs-bundles-equivariant-multiplicity-mirror-symmetry}


Fix $C$ a smooth complex projective curve of genus $g>1$.
Recall that a vector bundle $E$ on $C$ is called 
{\em stable} (resp. {\em semi-stable}) %if for all $\Phi$-invariant
if for all proper subbundles $F\subset E$ \beq \label{stable}{\deg(F)\over {\rm rank}(F)}
\source{very-stable-higgs-bundles-equivariant-multiplicity-mirror-symmetry}
< {\deg(E)\over {\rm rank}(E)}\eeq (resp. ${\deg(F)/ {\rm rank}(F)}
\leq {\deg(E)/ {\rm rank}(E)}$).

Similarly, a Higgs bundle $(E,\Phi)$ where $\Phi\in H^0(\End(E)\otimes K)$, is stable (resp. semi-stable) if \eqref{stable} holds for $\Phi$-invariant proper subbundles $F\subset E$. 

We fix integers $n\in \Z>0$ and $d\in \Z$.	The moduli space of semistable Higgs bundles $\M:=\M_{n}^{d}$ of rank $n$ and degree $d$ was constructed by \cite{hitchin} by gauge-theoretic and by \cite{nitsure} and \cite{simpson1} by algebraic geometric methods.  At its smooth points it carries a hyperk\"ahler metric, in particular an algebraic symplectic structure. It is  a normal \cite[Corollary 11.7]{simpson1b} quasi-projective variety. The open subset $\M^s\subset \M$ of stable points is precisely the smooth locus.  It is a coarse moduli space, but carries a universal Higgs bundle  \'etale locally on $\M^s$ \cite[Theorem 4.5.(4)]{simpson1}. 

For a Higgs bundle $(E,\Phi)$ we can compute the characteristic polynomial of the Higgs field as $\det(t-\Phi)=t^n+a_{1}t^{n-1}+\dots+a_n$, where $a_i\in H^0(C;K^i)$. This leads to the  map $$h:\M\to \oplus_{i=1}^n  H^0(C;K^i),$$ a proper, completely integrable Hamiltonian system \cite{hitchin-stable,simpson,nitsure} in particular, the fibres are Lagrangian at their smooth points. 

The moduli space $\M$ carries a $\C^\times$-action defined by $(E,\Phi)\mapsto (E,\lambda\Phi)$ for $\lambda\in \C^\times$. If we let $\C^\times$ act with weight $i$ on  $H^0(C;K^i)$ then $h$ becomes $\C^\times$-equivariant. As $h$ is proper and all the weights on $\oplus_{i=1}^n  H^0(C;K^i)$ are positive, we  deduce that $\M$ is semi-projective with respect to this $\T$-action. 

Finally, we recall that there is a canonical symplectic structure on $\M^s$ originally arising as part of its hyperk\"ahler structure. The tangent space at $(E,\Phi)\in \M^s$ can be identified (see e.g. \cite{biswas-ramanan})
\beq\label{tangent} T_{(E,\Phi)}\M^s\cong \H^1\left(C;\End(E)\stackrel{\ad(\Phi)}{\to} \End(E)\otimes K\right).\eeq Serre duality will give an isomorphism between this hypercohomology and its dual. This defines a symplectic form on $T_{(E,\Phi)}$. To see that it gives rise to a closed form $\omega$, we recall \cite[Theorem 4.3]{biswas-ramanan} that $\omega=d\theta$ where $$\theta:\H^1\left(C;\End(E)\stackrel{\ad(\Phi)}{\to}\End(E)\otimes K\right)\to \C$$ is given by first mapping to $H^1(C;\End(E))$ and then using  Serre duality to pair this element with $\Phi\in H^0(C;\End(E)\otimes K)$.  With respect to the $\T$-action we clearly have that $\lambda^*(\theta)=\lambda\theta$ and thus $\lambda^*(\omega)=\lambda\omega$ is also of weight $1$. 
\source{very-stable-higgs-bundles-equivariant-multiplicity-mirror-symmetry}


In this paper we will study the upward and downward flows of the Bialynicki-Birula partition on $\M$.  

\subsection{Fixed points of the $\T$-action $\M^{s\T}$}
First we recall a parametrisation of $\M^\T$. If $(E,\Phi)\in \M^{s\T}$ is a stable Higgs bundle fixed by the $\T$-action then we have for every $\lambda\in\T$ $$(E,\lambda\Phi)\cong (E,\Phi).$$ The isomorphism above gives a vector bundle automorphism $f_\lambda:E\to E$  making the diagram \beq\label{fixed}\begin{array}{ccc} E&\stackrel{\lambda \Phi}{\longrightarrow} & EK\\ \!\!\!\!\! f_\lambda \downarrow &&\,\,\,\,\,\downarrow f_\lambda \\ E &\stackrel{\Phi}{\longrightarrow}& EK \end{array}\eeq
\source{very-stable-higgs-bundles-equivariant-multiplicity-mirror-symmetry}
commutative. These define a fibrewise $\T$-action on $E$. Let $E=L_0\oplus \dots\oplus  L_k$ be the weight space decomposition of this $\T$-action, for subbundles $L_i<E$ and $0\leq k \leq n-1$. Let $\T$ act on $L_i$ by weight $w_i\in \Z$, in other words $f_\lambda$ acts on $L_i$ as multiplication by $\lambda^{w_i}$. %We can assume $w_1>w_2>\dots>w_k$. 
From the diagram \eqref{fixed} we see that $\Phi$ maps the weight $w_i$ space to the $w_i-1$ weight space. Using an overall scaling we can assume that $w_i=-i$. Then $\Phi(L_i)\subset L_{i+1}K$. 
We call the ordered partition $(\rank(L_0),...,\rank(L_k))$  of $n$ the {\em type} of the fixed point $(E,\Phi)\in \M^{s\T}$.  Similarly we can talk about the type of a component of the fixed point set  $F\in \pi_0(\M^{s\T})$ as the type is locally constant on $\M^{s\T}$. The latter can be seen by noting that \'etale locally we have a $\T$-equivariant universal bundle on $\M^s$ and on the connected component $F\in \pi_0(\M^{s\T})$  the weight space for a given weight of the $\T$-action forms a vector bundle. 

\begin{example}
\source{very-stable-higgs-bundles-equivariant-multiplicity-mirror-symmetry}
\notready For example when $E=L_0$ the type is $(n)$ and $\Phi=0$, thus $E$ is a semistable bundle. We will denote by $\calN$ the moduli space of rank $n$ degree $d$ semistable bundles which is embedded $\calN\subset \M^\T$, the component of $\M^\T$ of type $(n)$.
\end{example}

\begin{example}
\source{very-stable-higgs-bundles-equivariant-multiplicity-mirror-symmetry}
\notready	\label{ex111} At the other extreme are the fixed points of type $(1,\dots,1)$. Let $(E,\Phi)\in \M^{s\T}$ be a type $(1,...,1)$ fixed point of the $\T$-action. In other words $E=L_0\oplus\dots \oplus L_{n-1}$ is a direct sum of line bundles and the Higgs field $\Phi$ satisfies $\Phi(L_{i-1})\subset L_{i}K$. Such a Higgs bundle is determined by $L_0$ and the choice of non-zero  $$b_i:=\Phi|_{L_{i-1}}\in H^0(C;L_{i-1}^*L_{i}K).$$ Note that the isomorphism class of such a Higgs bundle only depends on the effective divisors $\delta_i:=\div(b_i)\in C^{[m_i]}$, where $m_i=\deg(L_{i-1}^*L_{i}K)$, and $L_0$. To simplify notation we will choose a divisor $\delta_0$ so that $L_0=\calO(\delta_0)$. We will write $\delta=(\delta_0,\delta_1,\dots,\delta_{n-1})$. We also denote by $\calE_{\d}$ the type $(1,\dots,1)$ Higgs bundle for which $$L_i:=(\delta_0+\delta_1+\dots+\delta_{i-1})K^{-i+1}$$ and  $$b_i:=s_{\delta_i}\in H^0(C;\calO(\delta_i))\cong H^0(C;L_{i-1}^*L_{i}K)$$ the defining section of $\calO(\delta_i)$. 
\source{very-stable-higgs-bundles-equivariant-multiplicity-mirror-symmetry}
	
	% We will denote this Higgs bundle by $\calE_{L_1,\bb}$, where $\bb:=(b_1,\dots,b_{n-1})$. We note that if $\bb^\prime:= (b^\prime_1,\dots,b^\prime_{n-1})$ gives another $\calE_{L_1,\bb^\prime}$ such that $\div(b_i)=\div(b_{i}^\prime)$ then we can use an automorphism of $L_1\oplus\dots \oplus L_n$ to show that $\calE_{L_1,\bb}\cong \calE_{L_1,\bb^\prime}$. We will also denote by $d_i:=\div(b_i)\in C^{m_i}$ a degree $m_i=\deg(b_i)$ effective divisor on $C$. For $\d:=(d_1,\dots,d_{n-1})\in  C^{[m_1]}\times \dots \times C^{[m_{n-1}]}$  we will denote by $\calE_{L_1,\d}$  the isomorphism class of the Higgs bundles $\calE_{L_1,\bb}$ above.
	
	Let $\ell:=\deg(L_0)=\deg(\delta_0)$ and  $m:=(m_1,\dots,m_{n-1})$ where $m_i=\deg(b_i)=\deg(\delta_i)$. Then the ambient component $\calE_{\d}\in F_{\mm}\in \pi_0(\M^\T)$ of the fixed point set $\M^\T$ is isomorphic to  \beq\label{type111fix}F_{\mm}\cong J_\ell(C)\times C^{[m_1]}\times \dots \times C^{[m_{n-1}]}.\eeq Here the Jacobian $J_\ell(C)$ means the moduli space of degree $\ell$ divisors on $C$ modulo principal ones, when we represent its elements by degree $\ell$ divisors $\delta_0$, or equivalently $J_\ell(C)$ can mean the moduli space of degree $\ell$ line bundles when we represent its elements by the line bundle $L_0$. 
\source{very-stable-higgs-bundles-equivariant-multiplicity-mirror-symmetry}
\end{example}
\begin{remark}
\source{very-stable-higgs-bundles-equivariant-multiplicity-mirror-symmetry}
\notready \label{remarkn2} For rank $n=2$ we have just two different types. The type $(2)$ fixed point component is $\calN\subset \M$ the moduli space of semistable bundles with zero Higgs field. The type $(1,1)$ components $F_{\ell,m_1}\cong J_\ell(C)\times C^{[m_1]}$ parametrise Higgs bundles of the form $E=L_0\oplus L_1$ and $$\Phi=\left(\begin{array}{cc}0&0\\b&0\end{array} \right),$$ where $L_0$ and $L_1$ are line bundles of degrees $\ell_0:=\ell$ and $\ell_1$ respectively, $b\in H^0(C;L_0^*L_1K)$ and $m_1=\deg(b)=\ell_1-\ell_0+2g-2$. For stability we need $m_1<2g-2$. When $d=\ell_0+\ell_1$ is fixed then $2\ell_0=d-m_1+(2g-2)$ and we will abbreviate  $F_{m_1}:=F_{\ell,m_1}$. We then have $g-1$ type $(1,1)$ fixed point components $F_{1}$, $F_3$, \dots, $F_{2g-3}$ when $d$ is odd and $g-1$ of them $F_0,\dots, F_{2g-4}$ when $d$ is even.  
\source{very-stable-higgs-bundles-equivariant-multiplicity-mirror-symmetry}
\end{remark}
\subsection{Upward flows on $\M$}

We  now describe the upward flows in the Bialynicki-Birula partition of $\M$. We know from Proposition~\ref{bbprop} that the upward flow $W^+_\calE$ for $\calE\in \M^\T$ is isomorphic to a vector space modelled on $T_\calE^+$. The following proposition describes which Higgs bundles belong to $W^+_\calE$. 

\begin{proposition}
\source{very-stable-higgs-bundles-equivariant-multiplicity-mirror-symmetry}
\notready\label{filtrationexists} Let $\calE=(E,\Phi)\in \M$ then \beq\label{limitdown}\lim_{\lambda\to 0} (E,\lambda\Phi)=(E^\prime,\Phi^\prime)\eeq for some   $(E^\prime,\Phi^\prime)\in \M^{s\T}$ if and only if the following hold \begin{enumerate} \item there exists a  filtration \beq\label{filt}0=E_0\subset E_1\subset \cdots E_{k-1}\subset E_k = E\eeq by subbundles such that \item for all $i$ \beq\label{compup}\Phi(E_i)\subset E_{i+1}K \eeq \item  and the induced maps $$gr_0(\Phi):E_i/E_{i+1}\to E_{i+1}/E_{i+2}K$$ satisfy \beq \label{limithiggs}(E_1/E_0\oplus E_2/E_1\oplus \cdots\oplus E_k/E_{k-1},gr_0(\Phi))\cong (gr(E),gr_0(\Phi))\cong(E^\prime,\Phi^\prime).\eeq \end{enumerate}
\source{very-stable-higgs-bundles-equivariant-multiplicity-mirror-symmetry}
	Furthermore the filtration \eqref{filt} with these properties is unique. 
\end{proposition}
\begin{proof} The proof is by a Higgs bundle analogue of the Rees construction. The vector bundle version is discussed in \cite{asok,heinloth,halpern}. 
	
	First let us assume that we have a filtration \eqref{filt} which is compatible with the Higgs field i.e. \eqref{compup} holds for all $i$. We denote by $E^\prime:=gr(E)$ and $\Phi^\prime=gr_0(\Phi)$ and assume that $(E^\prime,\Phi^\prime)\in \M^{s\T}$ is a stable Higgs bundle. 
	
	We can define a vector bundle $\tilde{E}$ over $\C\times C=\spec(\C[x])\times C$ given by the $\Z$-graded $\calO_C[x]$ -module \beq \label{zgrading}\tilde{E}:=\bigoplus_{i\in \Z} x^{-i} E_{-i}\eeq on $C$, where $x$ acts via the embedding $x^{-i} E_{-i}\subset x^{-i+1} E_{-i+1}$. Here $E_k=E$ for $k\geq n$ and $E_k=0$ for $k\leq 0$. Note, that we only have non-positive weights in the $\Z$-graded module \eqref{zgrading}.
\source{very-stable-higgs-bundles-equivariant-multiplicity-mirror-symmetry}
	
	We also get the vector bundle homomorphism $\tilde{\Phi}:\tilde{E}\to \tilde{E}\otimes K$ by $\tilde{\Phi}:x^i E_i\to x^{i+1} E_{i+1}\otimes K$ given by $\Phi|_{E_i}$. Furthermore the $\Z$-grading \eqref{zgrading} defines a $\T$-action on $\tilde{E}$ covering the standard action of $\T$ on $\C\cong \spec(\bigoplus_{i\in \Z_{\leq 0}} x^{-i})$. Under this $\T$-action $\lambda\in \T$  sends $\tilde{\Phi}$ to $\lambda^{-1}\tilde{\Phi}$. 
	This gives us a family of Higgs bundles over $C$ parametrized by $\C$. Over $\{1\}\times C$ we get the vector bundle $$\tilde{E}/(x-1)\tilde{E}\cong E$$ back on $C$, while $\tilde{\Phi}$  will induce precisely $\Phi$, yielding our original Higgs bundle $$(\tilde{E},\tilde{\Phi})_{\{1\}\times C}\cong (E,\Phi)$$ 
	The $\T$-equivariant structure on $(\tilde{E},\tilde{\Phi})$ shows that over $\{\lambda\}\times C$ for $0\neq\lambda\in\C$ we get \beq\label{etildelambda}(\tilde{E},\tilde{\Phi})_{\{\lambda\}\times C}\cong (E,\lambda\Phi)\eeq Over   $\{0\}\times C\in \C$ the vector bundle $\tilde{E}$ restricts as  $$\tilde{E}/(x)\tilde{E}\cong gr(E)=E^\prime$$ while $\tilde{\Phi}$ restricts as $gr_0(\Phi)=\Phi^\prime$, giving the Higgs bundle \beq\label{etildezero}(\tilde{E},\tilde{\Phi})_{\{0\}\times C}\cong (E^\prime,\Phi^\prime).\eeq
\source{very-stable-higgs-bundles-equivariant-multiplicity-mirror-symmetry}
	
	As $\M$ is a coarse moduli space we have a morphism $f:\C\to \M$ such that \beq \label{fdef}f(\lambda)=\left\{\begin{array}{ll} (E,\lambda\Phi)& \mbox{ when } \lambda\neq 0\\ (E^\prime,\Phi^\prime)&\mbox{ when } \lambda=0 \end{array}\right.\eeq which exactly means \eqref{limitdown}.
\source{very-stable-higgs-bundles-equivariant-multiplicity-mirror-symmetry}
	
	For the other direction we assume \eqref{limitdown} and note that as $(E^\prime,\Phi^\prime)$ is a stable Higgs bundle $(E,\lambda\Phi)$ is stable too, because stability is an open condition (c.f. \cite[Lemma 3.7]{simpson1} or \cite[Proposition 3.1]{nitsure}). 
	Let us assume the existence of a map like $f$. Then, as the obstruction in $H^2(\M^s,\T)$ vanishes on $f(\C)$, we get a family of stable Higgs bundles $(\tilde{E},\tilde{\Phi})$ over $\C\times C$ with the properties \eqref{etildelambda} and \eqref{etildezero}. As in \S\ref{equiuni} and \cite[\S 4]{hausel-thaddeus} we can construct a $\T$-equivariant structure on $\tilde{E}$ where $\lambda\in \T$ acts on $\tilde{\Phi}$ as $\lambda^{-1}\tilde{\Phi}$. In other words if $v\in \tilde{E}$ then $$\tilde{\Phi}(\lambda(v))=\lambda(\tilde{\Phi}(v)),$$ where the $\T$-action on $\tilde{E}K$ is induced from the $\T$-action on $\tilde{E}$ and the weight  one action on $K$, i.e. for $\omega\in K$ $\lambda(\omega):=\lambda \omega$; compare it with \eqref{univdiag}. Denote by \beq \label{filtdef}E_i=\{v\in E\cong \tilde{E}|_{\{1\}}| \lim_{\lambda\to 0} \lambda^i \lambda(v)  \mbox{ exists} \}.\eeq Note that \beq\label{limphi}\lambda^i\lambda(\Phi(v))=\lambda^i\lambda(\tilde{\Phi}|_{\{1\}}(v))=\lambda^i\tilde{\Phi}(\lambda(v))=\tilde{\Phi}(\lambda^i\lambda(v)).\eeq So if $v\in E_i$ then $\lim_{\lambda\to 0}\lambda^i\lambda(v)$ exists and so does $\lim_{\lambda\to 0}\lambda^i\lambda(\Phi(v))$. Recalling that $\lambda$ acts on $K$ with weight one, we see that $\Phi(v)\in E_{i+1}K$.  As the construction is locally trivial over $C$, we get an increasing filtration $E_i\subset E_{i+1}\subset E$ by subbundles of $E$. Furthermore $\Phi(E_i)\subset E_{i+1}K$ thus $\Phi$ is compatible with the filtration. For $\mu\in \T$ we compute $$\mu(\lim_{\lambda\to 0}\lambda^i\lambda(v)) =\lim_{\lambda\to 0}\lambda^i\mu(\lambda(v))=\mu^{-i}\lim_{\lambda\to 0}(\mu\lambda)^i(\mu\lambda)(v)=\mu^{-i}\lim_{\lambda\to 0}\lambda^i\lambda(v).$$ This way we get a map
\source{very-stable-higgs-bundles-equivariant-multiplicity-mirror-symmetry}
	
	 $$f_i:\begin{array}{ccc}E_i&\to &(\tilde{E}|_{\{0\}})_{-i}\\ v &\mapsto& \lim_{\lambda\to 0} \lambda^{i}\lambda(v) \end{array}$$ from $E_i$ to the weight $-i$ $\T$-isotypical component of the vector bundle with $\T$-action $\tilde{E}|_{\{0\}}$ on $C$. The kernel of this map consists of those $v\in E_i$ for which $$ \lim_{\lambda\to 0}\lambda^i\lambda(v)=0,$$ i.e  there exists a section  $s\in \Gamma(\C,\tilde{E}|_{\C\times \{\pi(v)\}}),$ where $\pi:E\to C$ is the projection of the vector bundle,  $s(0)=0$ and $s(\lambda)=\lambda^i\lambda(v)$. Then we can write $s=\lambda s^\prime$ for another $s^\prime\in \Gamma(\C,\tilde{E}|_{\C\times \{\pi(v)\}})$ and so $ \lim_{\lambda\to 0}\lambda^{i-1}\lambda(v)=s^\prime(0)$ exists showing $\ker(f_i)=E_{i-1}$. Thus $f_i$ induces $$f_i:E_i/E_{i-1}\xhookrightarrow{} (\tilde{E}|_{\{0\}})_i\subset \tilde{E}|_{\{0\}}, $$ and if we put them together we get $f:gr(E)\xhookrightarrow{} \tilde{E}|_{\{0\}}$ an injective bundle map between  vector bundles of the same rank, which thus must be an isomorphism. 
	
	Finally, we note that \eqref{limphi} implies that for $v\in E_i$ we have $$f_{i+1}(\Phi(v))=\tilde{\Phi}|_0(f_i(v))$$ thus $$(gr(E),gr_0(\Phi))\cong(\tilde{E}|_{\{0\}},\tilde{\Phi}|_{\{0\}})$$ completing the proof of existence of the required compatible filtration.
	
	In order to prove the uniqueness of the filtration we start with a compatible filtration \eqref{fullfiltration} as in the first part of the proof. Then we construct the $\T$-equivariant family of stable Higgs bundles $\tilde{E}$ over $\C\times C$ as in \eqref{zgrading}. Then we will show that the original filtration agrees with the one constructed in the second part of the proof with the definition \eqref{filtdef}. First we reformulate \eqref{filtdef} by saying that $E^\prime_{-i}\subset E$ is the subsheaf whose sections $s\in H^0(U,E^\prime_{-i})\subset H^0(U,E)$ on open subsets $U\subset C$ satisfy the condition that the section $\lambda^{-i}\lambda(s)\in H^0(\C^*\times U,\tilde{E})$ extends to a section in $H^0(\C\times U,\tilde{E})$. We note that the $\C[x,x^{-1}]$-module $H^0(\C^*\times U,\tilde{E})$ is the localization of the $\C[x]$-module $\bigoplus_{i\in \Z} x^{-i}E_{-i}$  at $x$, thus $$H^0(\C^*\times U,\tilde{E})\cong \bigoplus_{i\in \Z} x^{-i}\left(\bigoplus_{j\in \N} H^0(U;E_{-i+j})\right). $$ We see that the $\T$-invariant section $\lambda(s)\in H^0(\C^*\times U;\tilde{E})^\T$ agrees with $s=x^0s$ in this description  and so the section $\lambda^{-i}\lambda (s)\in H^0(\C^*\times U;\tilde{E}) $ agrees with $x^{-i}s$ in this description also. Finally, we see that $x^{-i}s$ extends to a section over $\C\times U$ if and only it is in the image of the restriction map $H^0(\C\times U, \tilde{E})\to H^0(\C^*\times U, \tilde{E})$. This happens exactly when $s\in H^0(U;E_{-i})$. Thus $E^\prime_{-i}=E_{-i}$. Therefore the filtration in \eqref{fullfiltration}
	satisfying \eqref{limithiggs} is unique.
	
	The proof of the Proposition is complete. 
\end{proof}

\begin{remark}
\source{very-stable-higgs-bundles-equivariant-multiplicity-mirror-symmetry}
\notready When $n=2$ and $(E,\Phi)\in \M$ then we can find the filtration in Proposition~\ref{filtrationexists}, giving the limiting Higgs bundle $\lim_{\lambda\to 0}(E,\lambda\Phi)\in \M^\T$ as follows. If $E$ is semistable then the limiting Higgs bundle is simply $(E,0)$. When $E$ is not stable then it has a unique (maximal) destabilizing subbundle $L\subset E$. The compatibility condition \eqref{compup} in this case is vacuous, and thus the limiting Higgs bundle is $\calE_{\delta_0,\delta_1}$ where $L=\calO(\delta_0)$, $\delta_1=\div(b)$ with $b:=\Phi|_L:L\to (E/L)\otimes K$. 
\end{remark}

\begin{remark}
\source{very-stable-higgs-bundles-equivariant-multiplicity-mirror-symmetry}
\notready\label{bbonpe} If we projectivise the total space of the $\T$-equivariant vector bundle $\tilde{E}$ as  ${\mathbb P}(\tilde{E}),$ then the map ${\mathbb P}(\tilde{E})\to \C$ will be proper, and $\T$-equivariant with respect to the induced $\T$-action on ${\mathbb P}(\tilde{E})$. If follows that this  $\T$-action on ${\mathbb P}(\tilde{E})$ is semi-projective. The fixed point set of the action is $$({\mathbb P}(\tilde{E}))^\T=\cup_i {\mathbb P}((\tilde{E}|_{\{0\}})_i).$$ The Bialynicki-Birula decomposition ${\mathbb P}(\tilde{E})=\cup_i W^+_i$ gives   $${\mathbb P}(E_i)\setminus {\mathbb P}(E_{i-1}) = W^+_{-i}\cap {\mathbb P}(\tilde{E})|_{\{1\}},$$ which follows from  the  proof of Proposition~\ref{filtrationexists}. Thus, amusingly, the filtration $E_i\subset E_{i+1}\subset E$ in Proposition~\ref{filtrationexists} describing the Bialynicki-Birula upward flows on $\M$, can be recovered from the Bialynicki-Birula decomposition on ${\mathbb P}(\tilde{E})$.
\source{very-stable-higgs-bundles-equivariant-multiplicity-mirror-symmetry}
	
	In fact, one can construct the Bialynicki-Birula partition on the total space of the projectivised equivariant universal bundle (which always exists over the stable locus, and extends over the whole of $\M$ using Simpson's framed moduli Higgs moduli space \cite[Theorem 4.10]{simpson1}, see \eqref{grassmann} for $k=1$) on $\M\times C$ which will contain the information on the filtration in Proposition~\ref{filtrationexists} which determines the limiting  Higgs bundle. 
\end{remark}

\begin{remark}
\source{very-stable-higgs-bundles-equivariant-multiplicity-mirror-symmetry}
\notready   In a differential geometric language Proposition~\ref{filtrationexists} appeared  in \cite[Proposition 4.2]{collier-wentworth}. Also a closely related partition of the de Rham moduli space of holomorphic connections $\M_{\rm DR}$ appeared in \cite{simpson}. In fact the algorithm in \cite{simpson} can be used verbatim to find the filtration in Proposition~\ref{filtrationexists} producing the limiting Higgs bundle. 
\end{remark}

\begin{remark}
\source{very-stable-higgs-bundles-equivariant-multiplicity-mirror-symmetry}
\notready \label{hitchinsection} Let us recall the construction of a section of the Hitchin map from \cite{hitchin-lie}.  Let ${{a}}=(a_1,\dots,a_n)\in H^0(C;K)\oplus \dots \oplus H^0(C;K^n)=\calA$ be a point in the Hitchin base. For a line bundle $L$ on $C$ let 
\source{very-stable-higgs-bundles-equivariant-multiplicity-mirror-symmetry}
	$$E_L:=L\oplus LK^{-1}\oplus\dots\oplus LK^{1-n}.$$ Take the Higgs field $\Phi_{{a}}:E_L\to E_LK$ given by the companion matrix: 	\beq\label{companion}\Phi_{{a}}:=\left(\begin{array}{ccccc}  0 &0& \dots &0& -a_n\\ 1 & 0 & \dots &0&-a_{n-1}\\ 0 & 1 & \dots &0&-a_{n-2}\\ \vdots & \vdots & \ddots &\vdots &\vdots \\ 0 & 0& \dots &1 & -a_1
\source{very-stable-higgs-bundles-equivariant-multiplicity-mirror-symmetry}
	\end{array}\right).\eeq  Denote the Higgs bundle $\calE_{L,{a}}:=(E_L,\Phi_{{a}})$. Define $E_i:=L\oplus LK^{-1}\oplus\dots\oplus LK^{-i+1}$. The filtration $0\subset E_1\dots \subset E_n=E_L$ satisfies $\Phi_{{a}}(E_i)\subset E_{i+1}K$ and the associated graded is $(gr(E_L),gr_0(\Phi_{ a}))=(E_L,\Phi_0)=\calE_{L,0}$. It is straightforward to check this is stable. We call $\calE_{L,0}$ a {\em uniformising Higgs bundle}\footnote{The terminology is motivated by the $n=2$ case when the $\PGL_2$ Higgs bundle associated to our uniformising Higgs bundle corresponds by  non-abelian Hodge theory to the representation $\pi_1(C)\to \PGL_2(\R)\subset \PGL_2(\C)$ giving the uniformising hyperbolic metric on $C$.}.   
	Now Proposition~\ref{filtrationexists} applies and shows that $\calE_{L,a}$ is stable too and $\lim_{\lambda\to 0} (E_L,\lambda \Phi_{ a})=(E_L,\Phi_0)=\calE_{L,0}$. Thus we have a map $s:\calA\to \M$ given by $a\mapsto \calE_a$. By construction  $h(\calE_{L,a})=a$ thus $s$ is a section of the Hitchin map $h:\M\to \calA$. Thus we see that $\calA\cong s(\calA)\subset W^+_{\calE_{L,0}}$ is a subvariety of an affine variety which is isomorphic to a vector space of the same 
	dimension, therefore $s(\calA)=W^+_{\calE_{L,0}}$, and so the upward flow $W^+_{\calE_{L,0}}$ is just a Hitchin section. 
	
	In particular, this implies that if $\calE\in W^+_{\calE_{L,0}}$, i.e. it has a full filtration as in Proposition~\ref{filtrationexists} such that the associated graded is isomorphic to $\calE_{L,0}$, then it has to be isomorphic to a Higgs bundle  of the form $(E_L,\Phi_a)$. For example, the underlying bundle must  split as a direct  sum of line bundles. 
	
	If we choose  $L=\calO_C$ then we have \beq\label{uniform}\calE_0:=\calE_{{\calO_C},0}=E_{\calO_C}\stackrel{\Phi_0}{\to}E_{\calO_C}\otimes K \eeq  the {\em canonical uniformising Higgs bundle}. This will give a {\em canonical section} $W^+_{\calE_0}\subset \M$ of the Hitchin map, corresponding to the structure sheaves of the spectral curves under  the BNR correspondence, see \eqref{canonicalBNR}. 
\source{very-stable-higgs-bundles-equivariant-multiplicity-mirror-symmetry}
	
	%	$ \calE_0=(E_0,\Phi_0)$ where $E_0:=K^{n-1}\oplus K^{n-2}\oplus\dots\oplus \calO$ and $\Phi_0|_{K^{i}}:K^i\to K^{i-1}K$ is given by the identity$$\Phi_0|_{K^{i}}=1\in H^0(C;Hom(K^{i}\to K^{i-1}K))\cong H^0(C;\calO)$$
	
	
\end{remark}

\subsection{Downward flows on $\M$}

Recall that the core of $\M$ is defined as $\calC:=\coprod_\alpha W^-_\alpha = \coprod_F W^-_F$. From Proposition~\ref{lagrangian} we know that $W^-_F\subset \M^s$ is a Lagrangian subvariety. Thus the core $\calC$ is Lagrangian at its smooth points. 

By recalling the $\T$-equivariant map $h:\M\to\calA$ where $\T$ acts on $\calA$ with positive weights we see that the core of $\calA$ is the origin $0\in \calA$, and so $\calC\subset h^{-1}(0)$. On the other hand $h$ is proper, therefore $h^{-1}(0)$ is projective, thus $h^{-1}(0)\subset \calC$. The zero fibre of the Hitchin map $h^{-1}(0)$ is called the {\em nilpotent cone}, because $h(E,\Phi)=0$ if and only if $\Phi$ is nilpotent. We will think of the nilpotent cone $N:=h^{-1}(0)$ as a subscheme  of $\M$. In particular, we will see later that most of its components are non-reduced. The argument above then implies 

\begin{proposition}
\source{very-stable-higgs-bundles-equivariant-multiplicity-mirror-symmetry}
\notready\label{nilpcore} The reduced scheme of the nilpotent cone coincides with the core: $N_{red}=\calC$. 
\source{very-stable-higgs-bundles-equivariant-multiplicity-mirror-symmetry}
\end{proposition}
\begin{remark}
\source{very-stable-higgs-bundles-equivariant-multiplicity-mirror-symmetry}
\notready
	The core is a Lagrangian subvariety. Thus Proposition~\ref{nilpcore} shows that the nilpotent cone is Lagrangian, which is a result of Laumon \cite[Theorem 3]{laumon}.
\end{remark}

We finish with the analogue of Proposition~\ref{filtrationexists} for the downward flow.

\begin{proposition}
\source{very-stable-higgs-bundles-equivariant-multiplicity-mirror-symmetry}
\notready\label{downwardfiltration} Let $\calE=(E,\Phi)\in \M$ then \bes\label{limitup}\lim_{\lambda \to \infty} (E,\lambda\Phi)=(E^\prime,\Phi^\prime)\ees for some   $(E^\prime,\Phi^\prime)\in \M^{s\T}$ if and only if the following hold 
\source{very-stable-higgs-bundles-equivariant-multiplicity-mirror-symmetry}
	\begin{enumerate} \item there exists a filtration \beq\label{filtdown}0=E_0\subset E_1\subset \cdots E_{k-1}\subset E_k = E\eeq by subbundles \item such that for all $i$ \bes\label{compdown}\Phi(E_{i+1})\subset E_{i}K\ees \item and the induced maps $gr_\infty(\Phi):E_i/E_{i-1}\to E_{i-1}/E_{i-2}K$ satisfy $$(E_1/E_0\oplus \cdots\oplus E_{k-1}/E_{k-2}\oplus  E_k/E_{k-1},gr_\infty(\Phi))\cong (gr(E),gr_\infty(\Phi))\cong(E^\prime,\Phi^\prime).$$
\source{very-stable-higgs-bundles-equivariant-multiplicity-mirror-symmetry}
		\end{enumerate}
	
	Additionally, the fibration \eqref{filtdown} with these properties is unique. 
\end{proposition}

\begin{proof} The proof is an appropriate modification of the proof of Proposition~\ref{filtrationexists}.  
	Just as in Remark~\ref{bbonpe} we can think of the filtration  \eqref{filtdown} on the Higgs bundle induced from the downward flows on the projectivised universal bundle. 
\end{proof}
